\documentclass[12pt]{article}
\usepackage{amsmath, amssymb}
\usepackage[margin=1in]{geometry}
\setlength{\parskip}{6pt}
\title{QUBO Formulation: Hamiltonian Path Problem}
\date{}
\begin{document}
\maketitle

\subsection*{Hamiltonian Path Problem}

\paragraph{Problem Statement.}
Given an undirected graph $G=(V,E)$ with $|V|=N$ vertices, the Hamiltonian Path Problem seeks a path that visits every vertex exactly once. We formulate this as an optimization problem where the goal is to find a vertex ordering that minimizes the number of consecutive pairs in the ordering that do not correspond to an edge in $G$. If a Hamiltonian path exists, the minimum value is zero; otherwise, the formulation identifies the ordering with the fewest missing edges.

\paragraph{Decision Variables.}
We introduce binary variables $x_{i,v} \in \{0,1\}$ for each position $i \in \{0, 1, \dots, N-1\}$ and each vertex $v \in V$. The variable $x_{i,v}$ equals $1$ if vertex $v$ is assigned to position $i$ in the path, and $0$ otherwise. A valid path corresponds to a permutation matrix, where each position is assigned exactly one vertex and each vertex is assigned to exactly one position.

\paragraph{QUBO Derivation.}
We construct a Quadratic Unconstrained Binary Optimization (QUBO) Hamiltonian $H(\mathbf{x})$ by combining an objective function with penalty terms that enforce the permutation constraints.

\textbf{Step 1: Objective Function.}
The objective is to minimize the number of missing consecutive edges. For a path defined by the variables, a consecutive edge exists between position $i$ and $i+1$ if the vertex at $i$ and the vertex at $i+1$ are connected by an edge in $E$. The term $(1 - x_{i,u} x_{i+1,v})$ equals $0$ if $x_{i,u}=x_{i+1,v}=1$ and $(u,v) \in E$, and $1$ otherwise. Summing over all positions and edges, the objective is:
\begin{align}
H_{\text{obj}} = \sum_{i=0}^{N-2} \sum_{(u,v) \in E} \left( 1 - x_{i,u} x_{i+1,v} \right).
\end{align}
Minimizing $H_{\text{obj}}$ minimizes the count of missing edges.

\textbf{Step 2: Constraints.}
We enforce two sets of equality constraints to ensure a valid permutation:
\begin{align}
\sum_{v \in V} x_{i,v} = 1, \quad \forall i \in \{0, \dots, N-1\} \quad \text{(Position uniqueness)}, \label{eq:pos_constraint} \\
\sum_{i=0}^{N-1} x_{i,v} = 1, \quad \forall v \in V \quad \text{(Vertex uniqueness)}. \label{eq:vert_constraint}
\end{align}

\textbf{Step 3: Penalty Terms.}
We convert the constraints to quadratic penalty terms using a weight $\lambda > 0$. For a constraint $\sum_k x_k = 1$, the squared penalty is $\lambda (\sum_k x_k - 1)^2$. Using the idempotent property $x_k^2 = x_k$ for binary variables, we expand:
\begin{align}
\lambda \left( \sum_k x_k - 1 \right)^2 &= \lambda \left( \sum_k x_k^2 + \sum_{k \neq l} x_k x_l - 2 \sum_k x_k + 1 \right) \nonumber \\
&= \lambda \left( \sum_k x_k + \sum_{k \neq l} x_k x_l - 2 \sum_k x_k + 1 \right) \nonumber \\
&= \lambda \left( \sum_{k \neq l} x_k x_l - \sum_k x_k + 1 \right).
\end{align}
Applying this to both constraints, the total penalty Hamiltonian is:
\begin{align}
H_{\text{pen}} &= \lambda \sum_{i=0}^{N-1} \left( \sum_{v \neq w} x_{i,v} x_{i,w} - \sum_{v \in V} x_{i,v} + 1 \right) \nonumber \\
&\quad + \lambda \sum_{v \in V} \left( \sum_{i \neq j} x_{i,v} x_{j,v} - \sum_{i=0}^{N-1} x_{i,v} + 1 \right).
\end{align}

\textbf{Step 4: Combined Hamiltonian.}
The total QUBO Hamiltonian $H(\mathbf{x}) = H_{\text{obj}} + H_{\text{pen}}$ is:
\begin{align}
H(\mathbf{x}) &= \sum_{i=0}^{N-2} \sum_{(u,v) \in E} \left( 1 - x_{i,u} x_{i+1,v} \right) \nonumber \\
&\quad + \lambda \sum_{i=0}^{N-1} \left( \sum_{v \neq w} x_{i,v} x_{i,w} - \sum_{v \in V} x_{i,v} + 1 \right) \nonumber \\
&\quad + \lambda \sum_{v \in V} \left( \sum_{i \neq j} x_{i,v} x_{j,v} - \sum_{i=0}^{N-1} x_{i,v} + 1 \right).
\end{align}

\textbf{Step 5: Q Matrix and Offset.}
Collecting terms, we identify the quadratic coefficients $Q$, linear coefficients, and constant offset $c$. The Hamiltonian takes the form $H(\mathbf{x}) = \mathbf{x}^T Q \mathbf{x} + c$. The non-zero entries of $Q$ are:
\begin{align}
Q_{(i,u),(i+1,v)} &= -1, \quad \forall i \in \{0, \dots, N-2\}, \, (u,v) \in E, \\
Q_{(i,v),(i,w)} &= \lambda, \quad \forall i \in \{0, \dots, N-1\}, \, v, w \in V, \, v \neq w, \\
Q_{(i,v),(j,v)} &= \lambda, \quad \forall v \in V, \, i, j \in \{0, \dots, N-1\}, \, i \neq j, \\
Q_{(i,v),(i,v)} &= -2\lambda, \quad \forall i \in \{0, \dots, N-1\}, \, v \in V.
\end{align}
All other entries are zero. The constant offset $c$ is the sum of constant terms from the objective and penalties:
\begin{align}
c = (N-1)|E| + 2\lambda N.
\end{align}

\paragraph{Penalty Weight Justification.}
The penalty weight $\lambda$ must be chosen such that any violation of the constraints incurs an energy cost strictly greater than the maximum possible gain from the objective function. The maximum magnitude of any coefficient in the objective function is $1$. A standard sufficient condition is $\lambda > \max |Q_{\text{obj}}|$. Setting $\lambda = 2$ satisfies $\lambda > 1$, ensuring that the penalty for violating a constraint dominates any potential reduction in the objective value. Thus, the global minimum of $H(\mathbf{x})$ over $\{0,1\}^{N^2}$ corresponds to a feasible solution that minimizes the original objective.

\paragraph{Correctness Argument.}
Minimizing $H(\mathbf{x})$ recovers the optimal Hamiltonian path (or best approximation) because: (i) for any feasible assignment satisfying the permutation constraints, the penalty terms vanish, and $H(\mathbf{x})$ reduces to the objective function; (ii) for any infeasible assignment, at least one constraint is violated, incurring a penalty of at least $\lambda$. Since $\lambda$ exceeds the maximum possible objective coefficient magnitude, the energy of any infeasible solution is strictly higher than that of the best feasible solution. Consequently, the global minimum must be feasible and optimal with respect to the original problem.

\paragraph{Illustrative Example.}
Consider an instance with $N=3$ vertices $V=\{0,1,2\}$ and a complete graph with self-loops, i.e., $E = V \times V$, so $|E|=9$. The variables are $x_{i,v}$ for $i,v \in \{0,1,2\}$, totaling $9$ variables. We set $\lambda = 2$.

The objective function becomes:
\begin{align}
H_{\text{obj}} = \sum_{i=0}^{1} \sum_{u=0}^{2} \sum_{v=0}^{2} \left( 1 - x_{i,u} x_{i+1,v} \right).
\end{align}
The penalty terms use $\lambda=2$. The constant offset is:
\begin{align}
c = (3-1) \times 9 + 2 \times 2 \times 3 = 18 + 12 = 30.
\end{align}
The QUBO Hamiltonian $H(\mathbf{x})$ includes quadratic terms $-x_{i,u} x_{i+1,v}$ for all $u,v$, penalty terms $2 x_{i,v} x_{i,w}$ for $v \neq w$, $2 x_{i,v} x_{j,v}$ for $i \neq j$, linear terms $-4 x_{i,v}$, and offset $30$.

For any feasible permutation, e.g., $x_{0,0}=x_{1,1}=x_{2,2}=1$ and all others $0$, the objective terms sum to $-\sum_{i=0}^{1} \sum_{u,v} 1 = -2$ (since $\sum_{u,v} x_{i,u} x_{i+1,v} = 1$ for each step). The penalties are zero. The total energy is:
\begin{align}
H(\mathbf{x}) = -2 + 30 = 28.
\end{align}
Any feasible solution yields energy $28$. Infeasible solutions have strictly higher energy due to the penalty weight $\lambda=2$. Thus, the minimum energy is $28$, achieved by any valid vertex ordering.

\end{document}